
\documentclass{article}
\usepackage{colortbl}
\usepackage{makecell}
\usepackage{multirow}
\usepackage{supertabular}

\begin{document}

\newcounter{utterance}

\centering \large Interaction Transcript for game `clean\_up', experiment `1\_medium\_5obj\_fr', episode 1 with Qwen3.6{-}35B{-}A3B{-}FP8{-}without{-}reasoning.
\vspace{24pt}

{ \footnotesize  \setcounter{utterance}{1}
\setlength{\tabcolsep}{0pt}
\begin{supertabular}{c@{$\;$}|p{.15\linewidth}@{}p{.15\linewidth}p{.15\linewidth}p{.15\linewidth}p{.15\linewidth}p{.15\linewidth}}
    \# & $\;$A & \multicolumn{4}{c}{Game Master} & $\;\:$B\\
    \hline

    \theutterance \stepcounter{utterance}  
    & & \multicolumn{4}{p{0.6\linewidth}}{
        \cellcolor[rgb]{0.9,0.9,0.9}{
            \makecell[{{p{\linewidth}}}]{
                \texttt{\tiny{[P1$\langle$GM]}}
                \texttt{Je suis votre maître de jeu, et vous jouez à un jeu collaboratif avec la grille suivante comme plateau de jeu :} \\
\\ 
\texttt{```} \\
\texttt{ 1234567} \\
\texttt{╔═══╤═╤═╗} \\
\texttt{║C◌◌│◌│◌║ 1} \\
\texttt{║◌I◌│◌├─╢ 2} \\
\texttt{║◌◌◌├─┘◌║ 3} \\
\texttt{║◌◌◌│◌◌◌║ 4} \\
\texttt{╟─┬─┼───╢ 5} \\
\texttt{║◌│◌│◌WT║ 6} \\
\texttt{║◌│◌│◌◌H║ 7} \\
\texttt{╚═╧═╧═══╝} \\
\\ 
\texttt{```} \\
\\ 
\texttt{* Le bord supérieur affiche des coordonnées x croissantes vers la droite, et le bord droit des coordonnées y croissantes vers le bas.} \\
\texttt{* Les objets suivants sont placés aléatoirement sur votre grille : 'W', 'I', 'T', 'C', 'H'.} \\
\\ 
\\ 
\texttt{L'autre joueur voit une variante du plateau de jeu, où les objets sont placés à des emplacements aléatoires différents. Vous ne pouvez pas voir le plateau de l'autre joueur et il ne peut pas voir le vôtre.} \\
\\ 
\\ 
\texttt{**Objectif :**} \\
\\ 
\texttt{Les deux joueurs doivent déplacer les objets sur leur arrière{-}plan respectif de manière à ce que des objets identiques se retrouvent aux mêmes coordonnées. Vous devez communiquer avec l'autre joueur pour convenir d'une configuration d'objectif commune.} \\
\\ 
\\ 
\texttt{**Règles :**} \\
\\ 
\texttt{* À chaque tour, vous pouvez envoyer exactement l'une des deux commandes suivantes :} \\
\texttt{ 1. `DIRE: <MESSAGE>` : pour envoyer un message (tout ce qui se trouve jusqu'au prochain saut de ligne) à l'autre joueur. Je le transmettrai à votre partenaire.} \\
\texttt{2. `DÉPLACER: <OBJET>, (<X>, <Y>)` : pour déplacer un objet vers une nouvelle position, où `<X>` est la colonne et `<Y>` est la ligne. Je vous informerai si votre déplacement a réussi ou non.} \\
\texttt{* Si vous ne respectez pas le format, ou si vous envoyez plusieurs commandes à la fois, je dois vous pénaliser.} \\
\texttt{* Si les deux joueurs accumulent plus de 12 pénalités, vous perdez tous les deux la partie.} \\
\texttt{* Il est essentiel que vous communiquiez avec l'autre joueur concernant votre état d'objectif ! La *seule* façon de transmettre votre stratégie à l'autre joueur est d'utiliser la commande `DIRE: <MESSAGE>` !} \\
\\ 
\\ 
\texttt{**Déplacement d'objets**} \\
\\ 
\texttt{* Vous ne pouvez déplacer des objets que vers des cellules situées à l'intérieur des limites de la grille. La cellule cible doit être vide, c'est{-}à{-}dire qu'elle ne doit contenir que le symbole '◌'.} \\
\texttt{* Si vous essayez de déplacer un objet vers un endroit qui n'est pas vide, ou si vous essayez de le déplacer en dehors de la grille, je dois vous pénaliser. Vous avez droit à un nouvel essai.} \\
\texttt{* Avant d'effectuer un déplacement, vérifiez bien que l'emplacement cible est vide et qu'il ne contient aucune lettre, aucun cadre ou aucune ligne !} \\
\\ 
\\ 
\texttt{**Fin de partie**} \\
\\ 
\texttt{Si vous pensez avoir atteint l'objectif d'aligner tous les objets, vous pouvez demander à l'autre joueur de terminer la partie en envoyant `DIRE: terminé ?`. Si l'autre joueur vous demande de terminer la partie et que vous répondez `DIRE: terminé !`, la partie se termine.} \\
\\ 
\texttt{Les deux joueurs gagnent si la partie se termine en 20 tours, un tour étant défini comme l'envoi d'une commande valide par chacun des deux joueurs.} \\
\\ 
\texttt{**Score :**} \\
\\ 
\texttt{Plus les objets identiques sont proches les uns des autres sur les deux plateaux de jeu, plus vous obtenez de points. Les pénalités réduisent vos points. Pouvez{-}vous battre le record ?} \\
\\ 
\\ 
\texttt{Envoyez un message à l'autre joueur pour commencer le jeu !} \\
            }
        }
    }
    & \\ \\

    \theutterance \stepcounter{utterance}  
    & \multicolumn{4}{p{0.6\linewidth}}{
        \cellcolor[rgb]{0.9,0.9,0.9}{
            \makecell[{{p{\linewidth}}}]{
                \texttt{\tiny{[P1$\rangle$GM]}}
                \texttt{DIRE: Bonjour ! Je vois les objets C, I, T, W, H sur ma grille. Voici leurs positions actuelles :} \\
\texttt{{-} C en (1, 1)} \\
\texttt{{-} I en (2, 2)} \\
\texttt{{-} T en (6, 6)} \\
\texttt{{-} W en (6, 5)} \\
\texttt{{-} H en (6, 7)} \\
\\ 
\texttt{Proposons{-}nous de viser une configuration cible commune. Par exemple, alignons tous les objets sur la diagonale principale ou sur une ligne spécifique pour simplifier. Quelle est votre position actuelle pour chaque objet ?} \\
            }
        }
    }
    & & \\ \\

    \theutterance \stepcounter{utterance}  
    & & & \multicolumn{2}{p{0.3\linewidth}}{
        \cellcolor[rgb]{0.9,0.9,0.9}{
            \makecell[{{p{\linewidth}}}]{
                \texttt{\tiny{[GM$|$GM]}}
                \texttt{DIRE: Bonjour ! Je vois les objets C, I, T, W, H sur ma grille. Voici leurs positions actuelles :} \\
\texttt{{-} C en (1, 1)} \\
\texttt{{-} I en (2, 2)} \\
\texttt{{-} T en (6, 6)} \\
\texttt{{-} W en (6, 5)} \\
\texttt{{-} H en (6, 7)} \\
\\ 
\texttt{Proposons{-}nous de viser une configuration cible commune. Par exemple, alignons tous les objets sur la diagonale principale ou sur une ligne spécifique pour simplifier. Quelle est votre position actuelle pour chaque objet ?} \\
            }
        }
    }
    & & \\ \\

    \theutterance \stepcounter{utterance}  
    & & & \multicolumn{2}{p{0.3\linewidth}}{
        \cellcolor[rgb]{0.9,0.9,0.9}{
            \makecell[{{p{\linewidth}}}]{
                \texttt{\tiny{[GM$|$GM]}}
                \texttt{Invalid format: tail is not empty: '} \\
\texttt{{-} C en (1, 1)} \\
\texttt{{-} I en (2, 2)} \\
\texttt{{-} T en (6, 6)} \\
\texttt{{-} W en (6, 5)} \\
\texttt{{-} H en (6, 7)} \\
\\ 
\texttt{Proposons{-}nous de viser une configuration cible commune. Par exemple, alignons tous les objets sur la diagonale principale ou sur une ligne spécifique pour simplifier. Quelle est votre position actuelle pour chaque objet ?} \\
            }
        }
    }
    & & \\ \\

    \theutterance \stepcounter{utterance}  
    & & \multicolumn{4}{p{0.6\linewidth}}{
        \cellcolor[rgb]{0.9,0.9,0.9}{
            \makecell[{{p{\linewidth}}}]{
                \texttt{\tiny{[P1$\langle$GM]}}
                \texttt{Pénalité : Votre message ne doit rien contenir après la commande !} \\
\\ 
\texttt{Assurez{-}vous que votre réponse ne contienne que `DIRE: <MESSAGE>` ou `DÉPLACER: <OBJET>, (<X>, <Y>)`, et rien d'autre !} \\
\\ 
\texttt{Vous avez collectivement accumulé 1 pénalités sur 12.} \\
\texttt{Veuillez réessayer !} \\
            }
        }
    }
    & \\ \\

    \theutterance \stepcounter{utterance}  
    & \multicolumn{4}{p{0.6\linewidth}}{
        \cellcolor[rgb]{0.9,0.9,0.9}{
            \makecell[{{p{\linewidth}}}]{
                \texttt{\tiny{[P1$\rangle$GM]}}
                \texttt{DIRE: Bonjour ! Je vois les objets C, I, T, W, H sur ma grille. Voici leurs positions actuelles : C en (1, 1), I en (2, 2), T en (6, 6), W en (6, 5), H en (6, 7). Proposons{-}nous de viser une configuration cible commune. Par exemple, alignons tous les objets sur la diagonale principale ou sur une ligne spécifique pour simplifier. Quelle est votre position actuelle pour chaque objet ?} \\
            }
        }
    }
    & & \\ \\

    \theutterance \stepcounter{utterance}  
    & & & \multicolumn{2}{p{0.3\linewidth}}{
        \cellcolor[rgb]{0.9,0.9,0.9}{
            \makecell[{{p{\linewidth}}}]{
                \texttt{\tiny{[GM$|$GM]}}
                \texttt{DIRE: Bonjour ! Je vois les objets C, I, T, W, H sur ma grille. Voici leurs positions actuelles : C en (1, 1), I en (2, 2), T en (6, 6), W en (6, 5), H en (6, 7). Proposons{-}nous de viser une configuration cible commune. Par exemple, alignons tous les objets sur la diagonale principale ou sur une ligne spécifique pour simplifier. Quelle est votre position actuelle pour chaque objet ?} \\
            }
        }
    }
    & & \\ \\

    \theutterance \stepcounter{utterance}  
    & & \multicolumn{4}{p{0.6\linewidth}}{
        \cellcolor[rgb]{0.9,0.9,0.9}{
            \makecell[{{p{\linewidth}}}]{
                \texttt{\tiny{[GM$\rangle$P2]}}
                \texttt{Je suis votre maître de jeu, et vous jouez à un jeu collaboratif avec la grille suivante comme plateau de jeu :} \\
\\ 
\texttt{```} \\
\texttt{ 1234567} \\
\texttt{╔═══╤═╤═╗} \\
\texttt{║◌◌◌│C│◌║ 1} \\
\texttt{║◌W◌│◌├─╢ 2} \\
\texttt{║◌◌◌├─┘◌║ 3} \\
\texttt{║◌◌◌│◌◌H║ 4} \\
\texttt{╟─┬─┼───╢ 5} \\
\texttt{║T│◌│I◌◌║ 6} \\
\texttt{║◌│◌│◌◌◌║ 7} \\
\texttt{╚═╧═╧═══╝} \\
\\ 
\texttt{```} \\
\\ 
\texttt{* Le bord supérieur affiche des coordonnées x croissantes vers la droite, et le bord droit des coordonnées y croissantes vers le bas.} \\
\texttt{* Les objets suivants sont placés aléatoirement sur votre grille : 'W', 'I', 'T', 'C', 'H'.} \\
\\ 
\\ 
\texttt{L'autre joueur voit une variante du plateau de jeu, où les objets sont placés à des emplacements aléatoires différents. Vous ne pouvez pas voir le plateau de l'autre joueur et il ne peut pas voir le vôtre.} \\
\\ 
\\ 
\texttt{**Objectif :**} \\
\\ 
\texttt{Les deux joueurs doivent déplacer les objets sur leur arrière{-}plan respectif de manière à ce que des objets identiques se retrouvent aux mêmes coordonnées. Vous devez communiquer avec l'autre joueur pour convenir d'une configuration d'objectif commune.} \\
\\ 
\\ 
\texttt{**Règles :**} \\
\\ 
\texttt{* À chaque tour, vous pouvez envoyer exactement l'une des deux commandes suivantes :} \\
\texttt{ 1. `DIRE: <MESSAGE>` : pour envoyer un message (tout ce qui se trouve jusqu'au prochain saut de ligne) à l'autre joueur. Je le transmettrai à votre partenaire.} \\
\texttt{2. `DÉPLACER: <OBJET>, (<X>, <Y>)` : pour déplacer un objet vers une nouvelle position, où `<X>` est la colonne et `<Y>` est la ligne. Je vous informerai si votre déplacement a réussi ou non.} \\
\texttt{* Si vous ne respectez pas le format, ou si vous envoyez plusieurs commandes à la fois, je dois vous pénaliser.} \\
\texttt{* Si les deux joueurs accumulent plus de 12 pénalités, vous perdez tous les deux la partie.} \\
\texttt{* Il est essentiel que vous communiquiez avec l'autre joueur concernant votre état d'objectif ! La *seule* façon de transmettre votre stratégie à l'autre joueur est d'utiliser la commande `DIRE: <MESSAGE>` !} \\
\\ 
\\ 
\texttt{**Déplacement d'objets**} \\
\\ 
\texttt{* Vous ne pouvez déplacer des objets que vers des cellules situées à l'intérieur des limites de la grille. La cellule cible doit être vide, c'est{-}à{-}dire qu'elle ne doit contenir que le symbole '◌'.} \\
\texttt{* Si vous essayez de déplacer un objet vers un endroit qui n'est pas vide, ou si vous essayez de le déplacer en dehors de la grille, je dois vous pénaliser. Vous avez droit à un nouvel essai.} \\
\texttt{* Avant d'effectuer un déplacement, vérifiez bien que l'emplacement cible est vide et qu'il ne contient aucune lettre, aucun cadre ou aucune ligne !} \\
\\ 
\\ 
\texttt{**Fin de partie**} \\
\\ 
\texttt{Si vous pensez avoir atteint l'objectif d'aligner tous les objets, vous pouvez demander à l'autre joueur de terminer la partie en envoyant `DIRE: terminé ?`. Si l'autre joueur vous demande de terminer la partie et que vous répondez `DIRE: terminé !`, la partie se termine.} \\
\\ 
\texttt{Les deux joueurs gagnent si la partie se termine en 20 tours, un tour étant défini comme l'envoi d'une commande valide par chacun des deux joueurs.} \\
\\ 
\texttt{**Score :**} \\
\\ 
\texttt{Plus les objets identiques sont proches les uns des autres sur les deux plateaux de jeu, plus vous obtenez de points. Les pénalités réduisent vos points. Pouvez{-}vous battre le record ?} \\
\\ 
\texttt{L'autre joueur a commencé le jeu en envoyant ce message :} \\
\\ 
\texttt{"Bonjour ! Je vois les objets C, I, T, W, H sur ma grille. Voici leurs positions actuelles : C en (1, 1), I en (2, 2), T en (6, 6), W en (6, 5), H en (6, 7). Proposons{-}nous de viser une configuration cible commune. Par exemple, alignons tous les objets sur la diagonale principale ou sur une ligne spécifique pour simplifier. Quelle est votre position actuelle pour chaque objet ?"} \\
\\ 
\texttt{Quelle est votre première commande ?} \\
            }
        }
    }
    & \\ \\

    \theutterance \stepcounter{utterance}  
    & & & \multicolumn{4}{p{0.6\linewidth}}{
        \cellcolor[rgb]{0.9,0.9,0.9}{
            \makecell[{{p{\linewidth}}}]{
                \texttt{\tiny{[GM$\langle$P2]}}
                \texttt{DIRE: Bonjour ! Je confirme avoir les objets W, C, H, T, I sur ma grille. Voici mes positions actuelles : W en (2, 2), C en (4, 1), H en (7, 4), T en (1, 6), I en (4, 6). Pour simplifier, je propose de viser la diagonale principale (coordonnées x=y). Cela semble être un bon point de départ commun. Êtes{-}vous d'accord pour essayer de placer tous les objets sur la diagonale ?} \\
            }
        }
    }
     \\ \\

    \theutterance \stepcounter{utterance}  
    & & & \multicolumn{2}{p{0.3\linewidth}}{
        \cellcolor[rgb]{0.9,0.9,0.9}{
            \makecell[{{p{\linewidth}}}]{
                \texttt{\tiny{[GM$|$GM]}}
                \texttt{DIRE: Bonjour ! Je confirme avoir les objets W, C, H, T, I sur ma grille. Voici mes positions actuelles : W en (2, 2), C en (4, 1), H en (7, 4), T en (1, 6), I en (4, 6). Pour simplifier, je propose de viser la diagonale principale (coordonnées x=y). Cela semble être un bon point de départ commun. Êtes{-}vous d'accord pour essayer de placer tous les objets sur la diagonale ?} \\
            }
        }
    }
    & & \\ \\

    \theutterance \stepcounter{utterance}  
    & & \multicolumn{4}{p{0.6\linewidth}}{
        \cellcolor[rgb]{0.9,0.9,0.9}{
            \makecell[{{p{\linewidth}}}]{
                \texttt{\tiny{[P1$\langle$GM]}}
                \texttt{Votre message a été transmis à l'autre joueur.} \\
\\ 
\\ 
\texttt{Vous jouez actuellement le tour 1 sur un maximum de 20.} \\
\texttt{Vous avez collectivement accumulé 1 pénalités sur 12.} \\
\texttt{L'autre joueur a envoyé ce message :} \\
\\ 
\\ 
\texttt{"Bonjour ! Je confirme avoir les objets W, C, H, T, I sur ma grille. Voici mes positions actuelles : W en (2, 2), C en (4, 1), H en (7, 4), T en (1, 6), I en (4, 6). Pour simplifier, je propose de viser la diagonale principale (coordonnées x=y). Cela semble être un bon point de départ commun. Êtes{-}vous d'accord pour essayer de placer tous les objets sur la diagonale ?"} \\
\\ 
\\ 
\texttt{Quelle est votre prochaine commande ?} \\
            }
        }
    }
    & \\ \\

    \theutterance \stepcounter{utterance}  
    & \multicolumn{4}{p{0.6\linewidth}}{
        \cellcolor[rgb]{0.9,0.9,0.9}{
            \makecell[{{p{\linewidth}}}]{
                \texttt{\tiny{[P1$\rangle$GM]}}
                \texttt{DIRE: C'est noté. Je suis d'accord pour viser la diagonale principale (x=y). Les positions cibles seront donc : C(1,1), I(2,2), T(3,3), W(4,4), H(5,5). Je vais commencer par déplacer I vers (2,2) qui est déjà sa position, donc je vais déplacer C vers (1,1) (déjà fait), puis je vais déplacer T vers (3,3).} \\
            }
        }
    }
    & & \\ \\

    \theutterance \stepcounter{utterance}  
    & & & \multicolumn{2}{p{0.3\linewidth}}{
        \cellcolor[rgb]{0.9,0.9,0.9}{
            \makecell[{{p{\linewidth}}}]{
                \texttt{\tiny{[GM$|$GM]}}
                \texttt{DIRE: C'est noté. Je suis d'accord pour viser la diagonale principale (x=y). Les positions cibles seront donc : C(1,1), I(2,2), T(3,3), W(4,4), H(5,5). Je vais commencer par déplacer I vers (2,2) qui est déjà sa position, donc je vais déplacer C vers (1,1) (déjà fait), puis je vais déplacer T vers (3,3).} \\
            }
        }
    }
    & & \\ \\

    \theutterance \stepcounter{utterance}  
    & & \multicolumn{4}{p{0.6\linewidth}}{
        \cellcolor[rgb]{0.9,0.9,0.9}{
            \makecell[{{p{\linewidth}}}]{
                \texttt{\tiny{[GM$\rangle$P2]}}
                \texttt{Votre message a été transmis à l'autre joueur.} \\
\\ 
\\ 
\texttt{Vous jouez actuellement le tour 2 sur un maximum de 20.} \\
\texttt{Vous avez collectivement accumulé 1 pénalités sur 12.} \\
\texttt{L'autre joueur a envoyé ce message :} \\
\\ 
\\ 
\texttt{"C'est noté. Je suis d'accord pour viser la diagonale principale (x=y). Les positions cibles seront donc : C(1,1), I(2,2), T(3,3), W(4,4), H(5,5). Je vais commencer par déplacer I vers (2,2) qui est déjà sa position, donc je vais déplacer C vers (1,1) (déjà fait), puis je vais déplacer T vers (3,3)."} \\
\\ 
\\ 
\texttt{Quelle est votre prochaine commande ?} \\
            }
        }
    }
    & \\ \\

    \theutterance \stepcounter{utterance}  
    & & & \multicolumn{4}{p{0.6\linewidth}}{
        \cellcolor[rgb]{0.9,0.9,0.9}{
            \makecell[{{p{\linewidth}}}]{
                \texttt{\tiny{[GM$\langle$P2]}}
                \texttt{DIRE: Parfait, la stratégie de la diagonale est validée. Cependant, il y a une petite confusion sur les positions cibles. Sur ma grille, W est déjà en (2,2), pas I. I est en (4,6). De plus, T est en (1,6) et H en (7,4). Pour la diagonale, nous devons adapter les cibles aux objets disponibles. Proposons cette correspondance : C{-}>(1,1), W{-}>(2,2), T{-}>(3,3), I{-}>(4,4), H{-}>(5,5). Je vais commencer par déplacer T de (1,6) vers (3,3).} \\
            }
        }
    }
     \\ \\

    \theutterance \stepcounter{utterance}  
    & & & \multicolumn{2}{p{0.3\linewidth}}{
        \cellcolor[rgb]{0.9,0.9,0.9}{
            \makecell[{{p{\linewidth}}}]{
                \texttt{\tiny{[GM$|$GM]}}
                \texttt{DIRE: Parfait, la stratégie de la diagonale est validée. Cependant, il y a une petite confusion sur les positions cibles. Sur ma grille, W est déjà en (2,2), pas I. I est en (4,6). De plus, T est en (1,6) et H en (7,4). Pour la diagonale, nous devons adapter les cibles aux objets disponibles. Proposons cette correspondance : C{-}>(1,1), W{-}>(2,2), T{-}>(3,3), I{-}>(4,4), H{-}>(5,5). Je vais commencer par déplacer T de (1,6) vers (3,3).} \\
            }
        }
    }
    & & \\ \\

    \theutterance \stepcounter{utterance}  
    & & \multicolumn{4}{p{0.6\linewidth}}{
        \cellcolor[rgb]{0.9,0.9,0.9}{
            \makecell[{{p{\linewidth}}}]{
                \texttt{\tiny{[P1$\langle$GM]}}
                \texttt{Votre message a été transmis à l'autre joueur.} \\
\\ 
\\ 
\texttt{Vous jouez actuellement le tour 2 sur un maximum de 20.} \\
\texttt{Vous avez collectivement accumulé 1 pénalités sur 12.} \\
\texttt{L'autre joueur a envoyé ce message :} \\
\\ 
\\ 
\texttt{"Parfait, la stratégie de la diagonale est validée. Cependant, il y a une petite confusion sur les positions cibles. Sur ma grille, W est déjà en (2,2), pas I. I est en (4,6). De plus, T est en (1,6) et H en (7,4). Pour la diagonale, nous devons adapter les cibles aux objets disponibles. Proposons cette correspondance : C{-}>(1,1), W{-}>(2,2), T{-}>(3,3), I{-}>(4,4), H{-}>(5,5). Je vais commencer par déplacer T de (1,6) vers (3,3)."} \\
\\ 
\\ 
\texttt{Quelle est votre prochaine commande ?} \\
            }
        }
    }
    & \\ \\

    \theutterance \stepcounter{utterance}  
    & \multicolumn{4}{p{0.6\linewidth}}{
        \cellcolor[rgb]{0.9,0.9,0.9}{
            \makecell[{{p{\linewidth}}}]{
                \texttt{\tiny{[P1$\rangle$GM]}}
                \texttt{DIRE: D'accord, je valide cette correspondance : C(1,1), W(2,2), T(3,3), I(4,4), H(5,5). Je vais déplacer T de (6,6) vers (3,3).} \\
            }
        }
    }
    & & \\ \\

    \theutterance \stepcounter{utterance}  
    & & & \multicolumn{2}{p{0.3\linewidth}}{
        \cellcolor[rgb]{0.9,0.9,0.9}{
            \makecell[{{p{\linewidth}}}]{
                \texttt{\tiny{[GM$|$GM]}}
                \texttt{DIRE: D'accord, je valide cette correspondance : C(1,1), W(2,2), T(3,3), I(4,4), H(5,5). Je vais déplacer T de (6,6) vers (3,3).} \\
            }
        }
    }
    & & \\ \\

    \theutterance \stepcounter{utterance}  
    & & \multicolumn{4}{p{0.6\linewidth}}{
        \cellcolor[rgb]{0.9,0.9,0.9}{
            \makecell[{{p{\linewidth}}}]{
                \texttt{\tiny{[GM$\rangle$P2]}}
                \texttt{Votre message a été transmis à l'autre joueur.} \\
\\ 
\\ 
\texttt{Vous jouez actuellement le tour 3 sur un maximum de 20.} \\
\texttt{Vous avez collectivement accumulé 1 pénalités sur 12.} \\
\texttt{L'autre joueur a envoyé ce message :} \\
\\ 
\\ 
\texttt{"D'accord, je valide cette correspondance : C(1,1), W(2,2), T(3,3), I(4,4), H(5,5). Je vais déplacer T de (6,6) vers (3,3)."} \\
\\ 
\\ 
\texttt{Quelle est votre prochaine commande ?} \\
            }
        }
    }
    & \\ \\

    \theutterance \stepcounter{utterance}  
    & & & \multicolumn{4}{p{0.6\linewidth}}{
        \cellcolor[rgb]{0.9,0.9,0.9}{
            \makecell[{{p{\linewidth}}}]{
                \texttt{\tiny{[GM$\langle$P2]}}
                \texttt{DÉPLACER: T, (3, 3)} \\
            }
        }
    }
     \\ \\

    \theutterance \stepcounter{utterance}  
    & & & \multicolumn{2}{p{0.3\linewidth}}{
        \cellcolor[rgb]{0.9,0.9,0.9}{
            \makecell[{{p{\linewidth}}}]{
                \texttt{\tiny{[GM$|$GM]}}
                \texttt{DÉPLACER: T, (3, 3)} \\
            }
        }
    }
    & & \\ \\

    \theutterance \stepcounter{utterance}  
    & & & \multicolumn{2}{p{0.3\linewidth}}{
        \cellcolor[rgb]{0.9,0.9,0.9}{
            \makecell[{{p{\linewidth}}}]{
                \texttt{\tiny{[GM$|$GM]}}
                \texttt{T moved to (3, 3)} \\
            }
        }
    }
    & & \\ \\

    \theutterance \stepcounter{utterance}  
    & & \multicolumn{4}{p{0.6\linewidth}}{
        \cellcolor[rgb]{0.9,0.9,0.9}{
            \makecell[{{p{\linewidth}}}]{
                \texttt{\tiny{[P1$\langle$GM]}}
                \texttt{Votre message a été transmis à l'autre joueur.} \\
\\ 
\\ 
\texttt{Vous jouez actuellement le tour 3 sur un maximum de 20.} \\
\texttt{Vous avez collectivement accumulé 1 pénalités sur 12.} \\
\texttt{L'autre joueur a déplacé un objet sur sa grille.} \\
\\ 
\\ 
\texttt{Quelle est votre prochaine commande ?} \\
            }
        }
    }
    & \\ \\

    \theutterance \stepcounter{utterance}  
    & \multicolumn{4}{p{0.6\linewidth}}{
        \cellcolor[rgb]{0.9,0.9,0.9}{
            \makecell[{{p{\linewidth}}}]{
                \texttt{\tiny{[P1$\rangle$GM]}}
                \texttt{DÉPLACER: T, (3, 3)} \\
            }
        }
    }
    & & \\ \\

    \theutterance \stepcounter{utterance}  
    & & & \multicolumn{2}{p{0.3\linewidth}}{
        \cellcolor[rgb]{0.9,0.9,0.9}{
            \makecell[{{p{\linewidth}}}]{
                \texttt{\tiny{[GM$|$GM]}}
                \texttt{DÉPLACER: T, (3, 3)} \\
            }
        }
    }
    & & \\ \\

    \theutterance \stepcounter{utterance}  
    & & & \multicolumn{2}{p{0.3\linewidth}}{
        \cellcolor[rgb]{0.9,0.9,0.9}{
            \makecell[{{p{\linewidth}}}]{
                \texttt{\tiny{[GM$|$GM]}}
                \texttt{T moved to (3, 3)} \\
            }
        }
    }
    & & \\ \\

    \theutterance \stepcounter{utterance}  
    & & \multicolumn{4}{p{0.6\linewidth}}{
        \cellcolor[rgb]{0.9,0.9,0.9}{
            \makecell[{{p{\linewidth}}}]{
                \texttt{\tiny{[GM$\rangle$P2]}}
                \texttt{Déplacement de 'T' vers (3,3) réussi. Votre grille mise à jour ressemble à ceci :} \\
\texttt{```} \\
\texttt{ 1234567} \\
\texttt{╔═══╤═╤═╗} \\
\texttt{║◌◌◌│C│◌║ 1} \\
\texttt{║◌W◌│◌├─╢ 2} \\
\texttt{║◌◌T├─┘◌║ 3} \\
\texttt{║◌◌◌│◌◌H║ 4} \\
\texttt{╟─┬─┼───╢ 5} \\
\texttt{║◌│◌│I◌◌║ 6} \\
\texttt{║◌│◌│◌◌◌║ 7} \\
\texttt{╚═╧═╧═══╝} \\
\texttt{```} \\
\texttt{Vous jouez actuellement le tour 4 sur un maximum de 20.} \\
\texttt{Vous avez collectivement accumulé 1 pénalités sur 12.} \\
\texttt{L'autre joueur a déplacé un objet sur sa grille.} \\
\\ 
\\ 
\texttt{Quelle est votre prochaine commande ?} \\
            }
        }
    }
    & \\ \\

    \theutterance \stepcounter{utterance}  
    & & & \multicolumn{4}{p{0.6\linewidth}}{
        \cellcolor[rgb]{0.9,0.9,0.9}{
            \makecell[{{p{\linewidth}}}]{
                \texttt{\tiny{[GM$\langle$P2]}}
                \texttt{DÉPLACER: I, (4, 4)} \\
            }
        }
    }
     \\ \\

    \theutterance \stepcounter{utterance}  
    & & & \multicolumn{2}{p{0.3\linewidth}}{
        \cellcolor[rgb]{0.9,0.9,0.9}{
            \makecell[{{p{\linewidth}}}]{
                \texttt{\tiny{[GM$|$GM]}}
                \texttt{DÉPLACER: I, (4, 4)} \\
            }
        }
    }
    & & \\ \\

    \theutterance \stepcounter{utterance}  
    & & & \multicolumn{2}{p{0.3\linewidth}}{
        \cellcolor[rgb]{0.9,0.9,0.9}{
            \makecell[{{p{\linewidth}}}]{
                \texttt{\tiny{[GM$|$GM]}}
                \texttt{Pénalité : (4,4) n'est pas vide, mais contient '│'.} \\
\\ 
\texttt{Vous avez collectivement accumulé 2 pénalités sur 12.} \\
\texttt{Veuillez réessayer !} \\
            }
        }
    }
    & & \\ \\

    \theutterance \stepcounter{utterance}  
    & & \multicolumn{4}{p{0.6\linewidth}}{
        \cellcolor[rgb]{0.9,0.9,0.9}{
            \makecell[{{p{\linewidth}}}]{
                \texttt{\tiny{[GM$\rangle$P2]}}
                \texttt{Pénalité : (4,4) n'est pas vide, mais contient '│'.} \\
\\ 
\texttt{Vous avez collectivement accumulé 2 pénalités sur 12.} \\
\texttt{Veuillez réessayer !} \\
            }
        }
    }
    & \\ \\

    \theutterance \stepcounter{utterance}  
    & & & \multicolumn{4}{p{0.6\linewidth}}{
        \cellcolor[rgb]{0.9,0.9,0.9}{
            \makecell[{{p{\linewidth}}}]{
                \texttt{\tiny{[GM$\langle$P2]}}
                \texttt{DÉPLACER: I, (4, 5)} \\
            }
        }
    }
     \\ \\

    \theutterance \stepcounter{utterance}  
    & & & \multicolumn{2}{p{0.3\linewidth}}{
        \cellcolor[rgb]{0.9,0.9,0.9}{
            \makecell[{{p{\linewidth}}}]{
                \texttt{\tiny{[GM$|$GM]}}
                \texttt{DÉPLACER: I, (4, 5)} \\
            }
        }
    }
    & & \\ \\

    \theutterance \stepcounter{utterance}  
    & & & \multicolumn{2}{p{0.3\linewidth}}{
        \cellcolor[rgb]{0.9,0.9,0.9}{
            \makecell[{{p{\linewidth}}}]{
                \texttt{\tiny{[GM$|$GM]}}
                \texttt{Pénalité : (4,5) n'est pas vide, mais contient '┼'.} \\
\\ 
\texttt{Vous avez collectivement accumulé 3 pénalités sur 12.} \\
\texttt{Veuillez réessayer !} \\
            }
        }
    }
    & & \\ \\

    \theutterance \stepcounter{utterance}  
    & & \multicolumn{4}{p{0.6\linewidth}}{
        \cellcolor[rgb]{0.9,0.9,0.9}{
            \makecell[{{p{\linewidth}}}]{
                \texttt{\tiny{[GM$\rangle$P2]}}
                \texttt{Pénalité : (4,5) n'est pas vide, mais contient '┼'.} \\
\\ 
\texttt{Vous avez collectivement accumulé 3 pénalités sur 12.} \\
\texttt{Veuillez réessayer !} \\
            }
        }
    }
    & \\ \\

    \theutterance \stepcounter{utterance}  
    & & & \multicolumn{4}{p{0.6\linewidth}}{
        \cellcolor[rgb]{0.9,0.9,0.9}{
            \makecell[{{p{\linewidth}}}]{
                \texttt{\tiny{[GM$\langle$P2]}}
                \texttt{DÉPLACER: I, (5, 4)} \\
            }
        }
    }
     \\ \\

    \theutterance \stepcounter{utterance}  
    & & & \multicolumn{2}{p{0.3\linewidth}}{
        \cellcolor[rgb]{0.9,0.9,0.9}{
            \makecell[{{p{\linewidth}}}]{
                \texttt{\tiny{[GM$|$GM]}}
                \texttt{DÉPLACER: I, (5, 4)} \\
            }
        }
    }
    & & \\ \\

    \theutterance \stepcounter{utterance}  
    & & & \multicolumn{2}{p{0.3\linewidth}}{
        \cellcolor[rgb]{0.9,0.9,0.9}{
            \makecell[{{p{\linewidth}}}]{
                \texttt{\tiny{[GM$|$GM]}}
                \texttt{I moved to (5, 4)} \\
            }
        }
    }
    & & \\ \\

    \theutterance \stepcounter{utterance}  
    & & \multicolumn{4}{p{0.6\linewidth}}{
        \cellcolor[rgb]{0.9,0.9,0.9}{
            \makecell[{{p{\linewidth}}}]{
                \texttt{\tiny{[P1$\langle$GM]}}
                \texttt{Déplacement de 'T' vers (3,3) réussi. Votre grille mise à jour ressemble à ceci :} \\
\texttt{```} \\
\texttt{ 1234567} \\
\texttt{╔═══╤═╤═╗} \\
\texttt{║C◌◌│◌│◌║ 1} \\
\texttt{║◌I◌│◌├─╢ 2} \\
\texttt{║◌◌T├─┘◌║ 3} \\
\texttt{║◌◌◌│◌◌◌║ 4} \\
\texttt{╟─┬─┼───╢ 5} \\
\texttt{║◌│◌│◌W◌║ 6} \\
\texttt{║◌│◌│◌◌H║ 7} \\
\texttt{╚═╧═╧═══╝} \\
\texttt{```} \\
\texttt{Vous jouez actuellement le tour 4 sur un maximum de 20.} \\
\texttt{Vous avez collectivement accumulé 3 pénalités sur 12.} \\
\texttt{L'autre joueur a déplacé un objet sur sa grille.} \\
\\ 
\\ 
\texttt{Quelle est votre prochaine commande ?} \\
            }
        }
    }
    & \\ \\

    \theutterance \stepcounter{utterance}  
    & \multicolumn{4}{p{0.6\linewidth}}{
        \cellcolor[rgb]{0.9,0.9,0.9}{
            \makecell[{{p{\linewidth}}}]{
                \texttt{\tiny{[P1$\rangle$GM]}}
                \texttt{DÉPLACER: I, (4, 4)} \\
            }
        }
    }
    & & \\ \\

    \theutterance \stepcounter{utterance}  
    & & & \multicolumn{2}{p{0.3\linewidth}}{
        \cellcolor[rgb]{0.9,0.9,0.9}{
            \makecell[{{p{\linewidth}}}]{
                \texttt{\tiny{[GM$|$GM]}}
                \texttt{DÉPLACER: I, (4, 4)} \\
            }
        }
    }
    & & \\ \\

    \theutterance \stepcounter{utterance}  
    & & & \multicolumn{2}{p{0.3\linewidth}}{
        \cellcolor[rgb]{0.9,0.9,0.9}{
            \makecell[{{p{\linewidth}}}]{
                \texttt{\tiny{[GM$|$GM]}}
                \texttt{Pénalité : (4,4) n'est pas vide, mais contient '│'.} \\
\\ 
\texttt{Vous avez collectivement accumulé 4 pénalités sur 12.} \\
\texttt{Veuillez réessayer !} \\
            }
        }
    }
    & & \\ \\

    \theutterance \stepcounter{utterance}  
    & & \multicolumn{4}{p{0.6\linewidth}}{
        \cellcolor[rgb]{0.9,0.9,0.9}{
            \makecell[{{p{\linewidth}}}]{
                \texttt{\tiny{[P1$\langle$GM]}}
                \texttt{Pénalité : (4,4) n'est pas vide, mais contient '│'.} \\
\\ 
\texttt{Vous avez collectivement accumulé 4 pénalités sur 12.} \\
\texttt{Veuillez réessayer !} \\
            }
        }
    }
    & \\ \\

    \theutterance \stepcounter{utterance}  
    & \multicolumn{4}{p{0.6\linewidth}}{
        \cellcolor[rgb]{0.9,0.9,0.9}{
            \makecell[{{p{\linewidth}}}]{
                \texttt{\tiny{[P1$\rangle$GM]}}
                \texttt{DÉPLACER: I, (4, 5)} \\
            }
        }
    }
    & & \\ \\

    \theutterance \stepcounter{utterance}  
    & & & \multicolumn{2}{p{0.3\linewidth}}{
        \cellcolor[rgb]{0.9,0.9,0.9}{
            \makecell[{{p{\linewidth}}}]{
                \texttt{\tiny{[GM$|$GM]}}
                \texttt{DÉPLACER: I, (4, 5)} \\
            }
        }
    }
    & & \\ \\

    \theutterance \stepcounter{utterance}  
    & & & \multicolumn{2}{p{0.3\linewidth}}{
        \cellcolor[rgb]{0.9,0.9,0.9}{
            \makecell[{{p{\linewidth}}}]{
                \texttt{\tiny{[GM$|$GM]}}
                \texttt{Pénalité : (4,5) n'est pas vide, mais contient '┼'.} \\
\\ 
\texttt{Vous avez collectivement accumulé 5 pénalités sur 12.} \\
\texttt{Veuillez réessayer !} \\
            }
        }
    }
    & & \\ \\

    \theutterance \stepcounter{utterance}  
    & & \multicolumn{4}{p{0.6\linewidth}}{
        \cellcolor[rgb]{0.9,0.9,0.9}{
            \makecell[{{p{\linewidth}}}]{
                \texttt{\tiny{[P1$\langle$GM]}}
                \texttt{Pénalité : (4,5) n'est pas vide, mais contient '┼'.} \\
\\ 
\texttt{Vous avez collectivement accumulé 5 pénalités sur 12.} \\
\texttt{Veuillez réessayer !} \\
            }
        }
    }
    & \\ \\

    \theutterance \stepcounter{utterance}  
    & \multicolumn{4}{p{0.6\linewidth}}{
        \cellcolor[rgb]{0.9,0.9,0.9}{
            \makecell[{{p{\linewidth}}}]{
                \texttt{\tiny{[P1$\rangle$GM]}}
                \texttt{DÉPLACER: I, (4, 3)} \\
            }
        }
    }
    & & \\ \\

    \theutterance \stepcounter{utterance}  
    & & & \multicolumn{2}{p{0.3\linewidth}}{
        \cellcolor[rgb]{0.9,0.9,0.9}{
            \makecell[{{p{\linewidth}}}]{
                \texttt{\tiny{[GM$|$GM]}}
                \texttt{DÉPLACER: I, (4, 3)} \\
            }
        }
    }
    & & \\ \\

    \theutterance \stepcounter{utterance}  
    & & & \multicolumn{2}{p{0.3\linewidth}}{
        \cellcolor[rgb]{0.9,0.9,0.9}{
            \makecell[{{p{\linewidth}}}]{
                \texttt{\tiny{[GM$|$GM]}}
                \texttt{Pénalité : (4,3) n'est pas vide, mais contient '├'.} \\
\\ 
\texttt{Vous avez collectivement accumulé 6 pénalités sur 12.} \\
\texttt{Veuillez réessayer !} \\
            }
        }
    }
    & & \\ \\

    \theutterance \stepcounter{utterance}  
    & & \multicolumn{4}{p{0.6\linewidth}}{
        \cellcolor[rgb]{0.9,0.9,0.9}{
            \makecell[{{p{\linewidth}}}]{
                \texttt{\tiny{[P1$\langle$GM]}}
                \texttt{Pénalité : (4,3) n'est pas vide, mais contient '├'.} \\
\\ 
\texttt{Vous avez collectivement accumulé 6 pénalités sur 12.} \\
\texttt{Veuillez réessayer !} \\
            }
        }
    }
    & \\ \\

    \theutterance \stepcounter{utterance}  
    & \multicolumn{4}{p{0.6\linewidth}}{
        \cellcolor[rgb]{0.9,0.9,0.9}{
            \makecell[{{p{\linewidth}}}]{
                \texttt{\tiny{[P1$\rangle$GM]}}
                \texttt{DÉPLACER: I, (4, 2)} \\
            }
        }
    }
    & & \\ \\

    \theutterance \stepcounter{utterance}  
    & & & \multicolumn{2}{p{0.3\linewidth}}{
        \cellcolor[rgb]{0.9,0.9,0.9}{
            \makecell[{{p{\linewidth}}}]{
                \texttt{\tiny{[GM$|$GM]}}
                \texttt{DÉPLACER: I, (4, 2)} \\
            }
        }
    }
    & & \\ \\

    \theutterance \stepcounter{utterance}  
    & & & \multicolumn{2}{p{0.3\linewidth}}{
        \cellcolor[rgb]{0.9,0.9,0.9}{
            \makecell[{{p{\linewidth}}}]{
                \texttt{\tiny{[GM$|$GM]}}
                \texttt{Pénalité : (4,2) n'est pas vide, mais contient '│'.} \\
\\ 
\texttt{Vous avez collectivement accumulé 7 pénalités sur 12.} \\
\texttt{Veuillez réessayer !} \\
            }
        }
    }
    & & \\ \\

    \theutterance \stepcounter{utterance}  
    & & \multicolumn{4}{p{0.6\linewidth}}{
        \cellcolor[rgb]{0.9,0.9,0.9}{
            \makecell[{{p{\linewidth}}}]{
                \texttt{\tiny{[P1$\langle$GM]}}
                \texttt{Pénalité : (4,2) n'est pas vide, mais contient '│'.} \\
\\ 
\texttt{Vous avez collectivement accumulé 7 pénalités sur 12.} \\
\texttt{Veuillez réessayer !} \\
            }
        }
    }
    & \\ \\

    \theutterance \stepcounter{utterance}  
    & \multicolumn{4}{p{0.6\linewidth}}{
        \cellcolor[rgb]{0.9,0.9,0.9}{
            \makecell[{{p{\linewidth}}}]{
                \texttt{\tiny{[P1$\rangle$GM]}}
                \texttt{DÉPLACER: I, (4, 1)} \\
            }
        }
    }
    & & \\ \\

    \theutterance \stepcounter{utterance}  
    & & & \multicolumn{2}{p{0.3\linewidth}}{
        \cellcolor[rgb]{0.9,0.9,0.9}{
            \makecell[{{p{\linewidth}}}]{
                \texttt{\tiny{[GM$|$GM]}}
                \texttt{DÉPLACER: I, (4, 1)} \\
            }
        }
    }
    & & \\ \\

    \theutterance \stepcounter{utterance}  
    & & & \multicolumn{2}{p{0.3\linewidth}}{
        \cellcolor[rgb]{0.9,0.9,0.9}{
            \makecell[{{p{\linewidth}}}]{
                \texttt{\tiny{[GM$|$GM]}}
                \texttt{Pénalité : (4,1) n'est pas vide, mais contient '│'.} \\
\\ 
\texttt{Vous avez collectivement accumulé 8 pénalités sur 12.} \\
\texttt{Veuillez réessayer !} \\
            }
        }
    }
    & & \\ \\

    \theutterance \stepcounter{utterance}  
    & & \multicolumn{4}{p{0.6\linewidth}}{
        \cellcolor[rgb]{0.9,0.9,0.9}{
            \makecell[{{p{\linewidth}}}]{
                \texttt{\tiny{[P1$\langle$GM]}}
                \texttt{Pénalité : (4,1) n'est pas vide, mais contient '│'.} \\
\\ 
\texttt{Vous avez collectivement accumulé 8 pénalités sur 12.} \\
\texttt{Veuillez réessayer !} \\
            }
        }
    }
    & \\ \\

    \theutterance \stepcounter{utterance}  
    & \multicolumn{4}{p{0.6\linewidth}}{
        \cellcolor[rgb]{0.9,0.9,0.9}{
            \makecell[{{p{\linewidth}}}]{
                \texttt{\tiny{[P1$\rangle$GM]}}
                \texttt{DÉPLACER: I, (5, 2)} \\
            }
        }
    }
    & & \\ \\

    \theutterance \stepcounter{utterance}  
    & & & \multicolumn{2}{p{0.3\linewidth}}{
        \cellcolor[rgb]{0.9,0.9,0.9}{
            \makecell[{{p{\linewidth}}}]{
                \texttt{\tiny{[GM$|$GM]}}
                \texttt{DÉPLACER: I, (5, 2)} \\
            }
        }
    }
    & & \\ \\

    \theutterance \stepcounter{utterance}  
    & & & \multicolumn{2}{p{0.3\linewidth}}{
        \cellcolor[rgb]{0.9,0.9,0.9}{
            \makecell[{{p{\linewidth}}}]{
                \texttt{\tiny{[GM$|$GM]}}
                \texttt{I moved to (5, 2)} \\
            }
        }
    }
    & & \\ \\

    \theutterance \stepcounter{utterance}  
    & & \multicolumn{4}{p{0.6\linewidth}}{
        \cellcolor[rgb]{0.9,0.9,0.9}{
            \makecell[{{p{\linewidth}}}]{
                \texttt{\tiny{[GM$\rangle$P2]}}
                \texttt{Déplacement de 'I' vers (5,4) réussi. Votre grille mise à jour ressemble à ceci :} \\
\texttt{```} \\
\texttt{ 1234567} \\
\texttt{╔═══╤═╤═╗} \\
\texttt{║◌◌◌│C│◌║ 1} \\
\texttt{║◌W◌│◌├─╢ 2} \\
\texttt{║◌◌T├─┘◌║ 3} \\
\texttt{║◌◌◌│I◌H║ 4} \\
\texttt{╟─┬─┼───╢ 5} \\
\texttt{║◌│◌│◌◌◌║ 6} \\
\texttt{║◌│◌│◌◌◌║ 7} \\
\texttt{╚═╧═╧═══╝} \\
\texttt{```} \\
\texttt{Vous jouez actuellement le tour 5 sur un maximum de 20.} \\
\texttt{Vous avez collectivement accumulé 8 pénalités sur 12.} \\
\texttt{L'autre joueur a déplacé un objet sur sa grille.} \\
\\ 
\\ 
\texttt{Quelle est votre prochaine commande ?} \\
            }
        }
    }
    & \\ \\

    \theutterance \stepcounter{utterance}  
    & & & \multicolumn{4}{p{0.6\linewidth}}{
        \cellcolor[rgb]{0.9,0.9,0.9}{
            \makecell[{{p{\linewidth}}}]{
                \texttt{\tiny{[GM$\langle$P2]}}
                \texttt{DÉPLACER: H, (5, 5)} \\
            }
        }
    }
     \\ \\

    \theutterance \stepcounter{utterance}  
    & & & \multicolumn{2}{p{0.3\linewidth}}{
        \cellcolor[rgb]{0.9,0.9,0.9}{
            \makecell[{{p{\linewidth}}}]{
                \texttt{\tiny{[GM$|$GM]}}
                \texttt{DÉPLACER: H, (5, 5)} \\
            }
        }
    }
    & & \\ \\

    \theutterance \stepcounter{utterance}  
    & & & \multicolumn{2}{p{0.3\linewidth}}{
        \cellcolor[rgb]{0.9,0.9,0.9}{
            \makecell[{{p{\linewidth}}}]{
                \texttt{\tiny{[GM$|$GM]}}
                \texttt{Pénalité : (5,5) n'est pas vide, mais contient '─'.} \\
\\ 
\texttt{Vous avez collectivement accumulé 9 pénalités sur 12.} \\
\texttt{Veuillez réessayer !} \\
            }
        }
    }
    & & \\ \\

    \theutterance \stepcounter{utterance}  
    & & \multicolumn{4}{p{0.6\linewidth}}{
        \cellcolor[rgb]{0.9,0.9,0.9}{
            \makecell[{{p{\linewidth}}}]{
                \texttt{\tiny{[GM$\rangle$P2]}}
                \texttt{Pénalité : (5,5) n'est pas vide, mais contient '─'.} \\
\\ 
\texttt{Vous avez collectivement accumulé 9 pénalités sur 12.} \\
\texttt{Veuillez réessayer !} \\
            }
        }
    }
    & \\ \\

    \theutterance \stepcounter{utterance}  
    & & & \multicolumn{4}{p{0.6\linewidth}}{
        \cellcolor[rgb]{0.9,0.9,0.9}{
            \makecell[{{p{\linewidth}}}]{
                \texttt{\tiny{[GM$\langle$P2]}}
                \texttt{DÉPLACER: H, (6, 5)} \\
            }
        }
    }
     \\ \\

    \theutterance \stepcounter{utterance}  
    & & & \multicolumn{2}{p{0.3\linewidth}}{
        \cellcolor[rgb]{0.9,0.9,0.9}{
            \makecell[{{p{\linewidth}}}]{
                \texttt{\tiny{[GM$|$GM]}}
                \texttt{DÉPLACER: H, (6, 5)} \\
            }
        }
    }
    & & \\ \\

    \theutterance \stepcounter{utterance}  
    & & & \multicolumn{2}{p{0.3\linewidth}}{
        \cellcolor[rgb]{0.9,0.9,0.9}{
            \makecell[{{p{\linewidth}}}]{
                \texttt{\tiny{[GM$|$GM]}}
                \texttt{Pénalité : (6,5) n'est pas vide, mais contient '─'.} \\
\\ 
\texttt{Vous avez collectivement accumulé 10 pénalités sur 12.} \\
\texttt{Veuillez réessayer !} \\
            }
        }
    }
    & & \\ \\

    \theutterance \stepcounter{utterance}  
    & & \multicolumn{4}{p{0.6\linewidth}}{
        \cellcolor[rgb]{0.9,0.9,0.9}{
            \makecell[{{p{\linewidth}}}]{
                \texttt{\tiny{[GM$\rangle$P2]}}
                \texttt{Pénalité : (6,5) n'est pas vide, mais contient '─'.} \\
\\ 
\texttt{Vous avez collectivement accumulé 10 pénalités sur 12.} \\
\texttt{Veuillez réessayer !} \\
            }
        }
    }
    & \\ \\

    \theutterance \stepcounter{utterance}  
    & & & \multicolumn{4}{p{0.6\linewidth}}{
        \cellcolor[rgb]{0.9,0.9,0.9}{
            \makecell[{{p{\linewidth}}}]{
                \texttt{\tiny{[GM$\langle$P2]}}
                \texttt{DÉPLACER: H, (5, 6)} \\
            }
        }
    }
     \\ \\

    \theutterance \stepcounter{utterance}  
    & & & \multicolumn{2}{p{0.3\linewidth}}{
        \cellcolor[rgb]{0.9,0.9,0.9}{
            \makecell[{{p{\linewidth}}}]{
                \texttt{\tiny{[GM$|$GM]}}
                \texttt{DÉPLACER: H, (5, 6)} \\
            }
        }
    }
    & & \\ \\

    \theutterance \stepcounter{utterance}  
    & & & \multicolumn{2}{p{0.3\linewidth}}{
        \cellcolor[rgb]{0.9,0.9,0.9}{
            \makecell[{{p{\linewidth}}}]{
                \texttt{\tiny{[GM$|$GM]}}
                \texttt{H moved to (5, 6)} \\
            }
        }
    }
    & & \\ \\

    \theutterance \stepcounter{utterance}  
    & & \multicolumn{4}{p{0.6\linewidth}}{
        \cellcolor[rgb]{0.9,0.9,0.9}{
            \makecell[{{p{\linewidth}}}]{
                \texttt{\tiny{[P1$\langle$GM]}}
                \texttt{Déplacement de 'I' vers (5,2) réussi. Votre grille mise à jour ressemble à ceci :} \\
\texttt{```} \\
\texttt{ 1234567} \\
\texttt{╔═══╤═╤═╗} \\
\texttt{║C◌◌│◌│◌║ 1} \\
\texttt{║◌◌◌│I├─╢ 2} \\
\texttt{║◌◌T├─┘◌║ 3} \\
\texttt{║◌◌◌│◌◌◌║ 4} \\
\texttt{╟─┬─┼───╢ 5} \\
\texttt{║◌│◌│◌W◌║ 6} \\
\texttt{║◌│◌│◌◌H║ 7} \\
\texttt{╚═╧═╧═══╝} \\
\texttt{```} \\
\texttt{Vous jouez actuellement le tour 5 sur un maximum de 20.} \\
\texttt{Vous avez collectivement accumulé 10 pénalités sur 12.} \\
\texttt{L'autre joueur a déplacé un objet sur sa grille.} \\
\\ 
\\ 
\texttt{Quelle est votre prochaine commande ?} \\
            }
        }
    }
    & \\ \\

    \theutterance \stepcounter{utterance}  
    & \multicolumn{4}{p{0.6\linewidth}}{
        \cellcolor[rgb]{0.9,0.9,0.9}{
            \makecell[{{p{\linewidth}}}]{
                \texttt{\tiny{[P1$\rangle$GM]}}
                \texttt{DIRE: Attention, nous avons beaucoup de pénalités. Je suis en (5,2) pour I. Je dois aller en (4,4). Le chemin est bloqué par les lignes verticales en x=4. Je vais essayer de contourner. Je vais déplacer I vers (5,4) d'abord.} \\
            }
        }
    }
    & & \\ \\

    \theutterance \stepcounter{utterance}  
    & & & \multicolumn{2}{p{0.3\linewidth}}{
        \cellcolor[rgb]{0.9,0.9,0.9}{
            \makecell[{{p{\linewidth}}}]{
                \texttt{\tiny{[GM$|$GM]}}
                \texttt{DIRE: Attention, nous avons beaucoup de pénalités. Je suis en (5,2) pour I. Je dois aller en (4,4). Le chemin est bloqué par les lignes verticales en x=4. Je vais essayer de contourner. Je vais déplacer I vers (5,4) d'abord.} \\
            }
        }
    }
    & & \\ \\

    \theutterance \stepcounter{utterance}  
    & & \multicolumn{4}{p{0.6\linewidth}}{
        \cellcolor[rgb]{0.9,0.9,0.9}{
            \makecell[{{p{\linewidth}}}]{
                \texttt{\tiny{[GM$\rangle$P2]}}
                \texttt{Déplacement de 'H' vers (5,6) réussi. Votre grille mise à jour ressemble à ceci :} \\
\texttt{```} \\
\texttt{ 1234567} \\
\texttt{╔═══╤═╤═╗} \\
\texttt{║◌◌◌│C│◌║ 1} \\
\texttt{║◌W◌│◌├─╢ 2} \\
\texttt{║◌◌T├─┘◌║ 3} \\
\texttt{║◌◌◌│I◌◌║ 4} \\
\texttt{╟─┬─┼───╢ 5} \\
\texttt{║◌│◌│H◌◌║ 6} \\
\texttt{║◌│◌│◌◌◌║ 7} \\
\texttt{╚═╧═╧═══╝} \\
\texttt{```} \\
\texttt{Vous jouez actuellement le tour 6 sur un maximum de 20.} \\
\texttt{Vous avez collectivement accumulé 10 pénalités sur 12.} \\
\texttt{L'autre joueur a envoyé ce message :} \\
\\ 
\\ 
\texttt{"Attention, nous avons beaucoup de pénalités. Je suis en (5,2) pour I. Je dois aller en (4,4). Le chemin est bloqué par les lignes verticales en x=4. Je vais essayer de contourner. Je vais déplacer I vers (5,4) d'abord."} \\
\\ 
\\ 
\texttt{Quelle est votre prochaine commande ?} \\
            }
        }
    }
    & \\ \\

    \theutterance \stepcounter{utterance}  
    & & & \multicolumn{4}{p{0.6\linewidth}}{
        \cellcolor[rgb]{0.9,0.9,0.9}{
            \makecell[{{p{\linewidth}}}]{
                \texttt{\tiny{[GM$\langle$P2]}}
                \texttt{DIRE: Attention, nous sommes à 10 pénalités. Il faut être très prudent. Pour moi, il reste à déplacer C en (1,1) (déjà fait), W en (2,2) (déjà fait), T en (3,3) (déjà fait), I en (4,4) et H en (5,5) (déplacé en 5,6). Je dois donc déplacer I vers (4,4). Vérifions si (4,4) est libre. Sur ma grille actuelle, (4,4) contient 'I' ? Non, I est en (5,4). Donc (4,4) est vide ? Regardons la grille : ligne 4, colonne 4. C'est 'I' ? Non, I est en (5,4). La case (4,4) est vide ? Oui, elle semble vide. Je vais essayer de déplacer I vers (4,4).} \\
            }
        }
    }
     \\ \\

    \theutterance \stepcounter{utterance}  
    & & & \multicolumn{2}{p{0.3\linewidth}}{
        \cellcolor[rgb]{0.9,0.9,0.9}{
            \makecell[{{p{\linewidth}}}]{
                \texttt{\tiny{[GM$|$GM]}}
                \texttt{DIRE: Attention, nous sommes à 10 pénalités. Il faut être très prudent. Pour moi, il reste à déplacer C en (1,1) (déjà fait), W en (2,2) (déjà fait), T en (3,3) (déjà fait), I en (4,4) et H en (5,5) (déplacé en 5,6). Je dois donc déplacer I vers (4,4). Vérifions si (4,4) est libre. Sur ma grille actuelle, (4,4) contient 'I' ? Non, I est en (5,4). Donc (4,4) est vide ? Regardons la grille : ligne 4, colonne 4. C'est 'I' ? Non, I est en (5,4). La case (4,4) est vide ? Oui, elle semble vide. Je vais essayer de déplacer I vers (4,4).} \\
            }
        }
    }
    & & \\ \\

    \theutterance \stepcounter{utterance}  
    & & \multicolumn{4}{p{0.6\linewidth}}{
        \cellcolor[rgb]{0.9,0.9,0.9}{
            \makecell[{{p{\linewidth}}}]{
                \texttt{\tiny{[P1$\langle$GM]}}
                \texttt{Votre message a été transmis à l'autre joueur.} \\
\\ 
\\ 
\texttt{Vous jouez actuellement le tour 6 sur un maximum de 20.} \\
\texttt{Vous avez collectivement accumulé 10 pénalités sur 12.} \\
\texttt{L'autre joueur a envoyé ce message :} \\
\\ 
\\ 
\texttt{"Attention, nous sommes à 10 pénalités. Il faut être très prudent. Pour moi, il reste à déplacer C en (1,1) (déjà fait), W en (2,2) (déjà fait), T en (3,3) (déjà fait), I en (4,4) et H en (5,5) (déplacé en 5,6). Je dois donc déplacer I vers (4,4). Vérifions si (4,4) est libre. Sur ma grille actuelle, (4,4) contient 'I' ? Non, I est en (5,4). Donc (4,4) est vide ? Regardons la grille : ligne 4, colonne 4. C'est 'I' ? Non, I est en (5,4). La case (4,4) est vide ? Oui, elle semble vide. Je vais essayer de déplacer I vers (4,4)."} \\
\\ 
\\ 
\texttt{Quelle est votre prochaine commande ?} \\
            }
        }
    }
    & \\ \\

    \theutterance \stepcounter{utterance}  
    & \multicolumn{4}{p{0.6\linewidth}}{
        \cellcolor[rgb]{0.9,0.9,0.9}{
            \makecell[{{p{\linewidth}}}]{
                \texttt{\tiny{[P1$\rangle$GM]}}
                \texttt{DÉPLACER: I, (4, 4)} \\
            }
        }
    }
    & & \\ \\

    \theutterance \stepcounter{utterance}  
    & & & \multicolumn{2}{p{0.3\linewidth}}{
        \cellcolor[rgb]{0.9,0.9,0.9}{
            \makecell[{{p{\linewidth}}}]{
                \texttt{\tiny{[GM$|$GM]}}
                \texttt{DÉPLACER: I, (4, 4)} \\
            }
        }
    }
    & & \\ \\

    \theutterance \stepcounter{utterance}  
    & & & \multicolumn{2}{p{0.3\linewidth}}{
        \cellcolor[rgb]{0.9,0.9,0.9}{
            \makecell[{{p{\linewidth}}}]{
                \texttt{\tiny{[GM$|$GM]}}
                \texttt{Pénalité : (4,4) n'est pas vide, mais contient '│'.} \\
\\ 
\texttt{Vous avez collectivement accumulé 11 pénalités sur 12.} \\
\texttt{Veuillez réessayer !} \\
            }
        }
    }
    & & \\ \\

    \theutterance \stepcounter{utterance}  
    & & \multicolumn{4}{p{0.6\linewidth}}{
        \cellcolor[rgb]{0.9,0.9,0.9}{
            \makecell[{{p{\linewidth}}}]{
                \texttt{\tiny{[P1$\langle$GM]}}
                \texttt{Pénalité : (4,4) n'est pas vide, mais contient '│'.} \\
\\ 
\texttt{Vous avez collectivement accumulé 11 pénalités sur 12.} \\
\texttt{Veuillez réessayer !} \\
            }
        }
    }
    & \\ \\

    \theutterance \stepcounter{utterance}  
    & \multicolumn{4}{p{0.6\linewidth}}{
        \cellcolor[rgb]{0.9,0.9,0.9}{
            \makecell[{{p{\linewidth}}}]{
                \texttt{\tiny{[P1$\rangle$GM]}}
                \texttt{DÉPLACER: I, (5, 4)} \\
            }
        }
    }
    & & \\ \\

    \theutterance \stepcounter{utterance}  
    & & & \multicolumn{2}{p{0.3\linewidth}}{
        \cellcolor[rgb]{0.9,0.9,0.9}{
            \makecell[{{p{\linewidth}}}]{
                \texttt{\tiny{[GM$|$GM]}}
                \texttt{DÉPLACER: I, (5, 4)} \\
            }
        }
    }
    & & \\ \\

    \theutterance \stepcounter{utterance}  
    & & & \multicolumn{2}{p{0.3\linewidth}}{
        \cellcolor[rgb]{0.9,0.9,0.9}{
            \makecell[{{p{\linewidth}}}]{
                \texttt{\tiny{[GM$|$GM]}}
                \texttt{I moved to (5, 4)} \\
            }
        }
    }
    & & \\ \\

    \theutterance \stepcounter{utterance}  
    & & \multicolumn{4}{p{0.6\linewidth}}{
        \cellcolor[rgb]{0.9,0.9,0.9}{
            \makecell[{{p{\linewidth}}}]{
                \texttt{\tiny{[GM$\rangle$P2]}}
                \texttt{Votre message a été transmis à l'autre joueur.} \\
\\ 
\\ 
\texttt{Vous jouez actuellement le tour 7 sur un maximum de 20.} \\
\texttt{Vous avez collectivement accumulé 11 pénalités sur 12.} \\
\texttt{L'autre joueur a déplacé un objet sur sa grille.} \\
\\ 
\\ 
\texttt{Quelle est votre prochaine commande ?} \\
            }
        }
    }
    & \\ \\

    \theutterance \stepcounter{utterance}  
    & & & \multicolumn{4}{p{0.6\linewidth}}{
        \cellcolor[rgb]{0.9,0.9,0.9}{
            \makecell[{{p{\linewidth}}}]{
                \texttt{\tiny{[GM$\langle$P2]}}
                \texttt{DÉPLACER: I, (4, 4)} \\
            }
        }
    }
     \\ \\

    \theutterance \stepcounter{utterance}  
    & & & \multicolumn{2}{p{0.3\linewidth}}{
        \cellcolor[rgb]{0.9,0.9,0.9}{
            \makecell[{{p{\linewidth}}}]{
                \texttt{\tiny{[GM$|$GM]}}
                \texttt{DÉPLACER: I, (4, 4)} \\
            }
        }
    }
    & & \\ \\

    \theutterance \stepcounter{utterance}  
    & & & \multicolumn{2}{p{0.3\linewidth}}{
        \cellcolor[rgb]{0.9,0.9,0.9}{
            \makecell[{{p{\linewidth}}}]{
                \texttt{\tiny{[GM$|$GM]}}
                \texttt{Pénalité : (4,4) n'est pas vide, mais contient '│'.} \\
\\ 
\texttt{Vous avez collectivement accumulé 12 pénalités sur 12.} \\
\texttt{Veuillez réessayer !} \\
            }
        }
    }
    & & \\ \\

    \theutterance \stepcounter{utterance}  
    & & \multicolumn{4}{p{0.6\linewidth}}{
        \cellcolor[rgb]{0.9,0.9,0.9}{
            \makecell[{{p{\linewidth}}}]{
                \texttt{\tiny{[GM$\rangle$P2]}}
                \texttt{Pénalité : (4,4) n'est pas vide, mais contient '│'.} \\
\\ 
\texttt{Vous avez collectivement accumulé 12 pénalités sur 12.} \\
\texttt{Veuillez réessayer !} \\
            }
        }
    }
    & \\ \\

    \theutterance \stepcounter{utterance}  
    & & & \multicolumn{4}{p{0.6\linewidth}}{
        \cellcolor[rgb]{0.9,0.9,0.9}{
            \makecell[{{p{\linewidth}}}]{
                \texttt{\tiny{[GM$\langle$P2]}}
                \texttt{DIRE: terminé ?} \\
            }
        }
    }
     \\ \\

    \theutterance \stepcounter{utterance}  
    & & & \multicolumn{2}{p{0.3\linewidth}}{
        \cellcolor[rgb]{0.9,0.9,0.9}{
            \makecell[{{p{\linewidth}}}]{
                \texttt{\tiny{[GM$|$GM]}}
                \texttt{DIRE: terminé ?} \\
            }
        }
    }
    & & \\ \\

    \theutterance \stepcounter{utterance}  
    & & \multicolumn{4}{p{0.6\linewidth}}{
        \cellcolor[rgb]{0.9,0.9,0.9}{
            \makecell[{{p{\linewidth}}}]{
                \texttt{\tiny{[P1$\langle$GM]}}
                \texttt{Déplacement de 'I' vers (5,4) réussi. Votre grille mise à jour ressemble à ceci :} \\
\texttt{```} \\
\texttt{ 1234567} \\
\texttt{╔═══╤═╤═╗} \\
\texttt{║C◌◌│◌│◌║ 1} \\
\texttt{║◌◌◌│◌├─╢ 2} \\
\texttt{║◌◌T├─┘◌║ 3} \\
\texttt{║◌◌◌│I◌◌║ 4} \\
\texttt{╟─┬─┼───╢ 5} \\
\texttt{║◌│◌│◌W◌║ 6} \\
\texttt{║◌│◌│◌◌H║ 7} \\
\texttt{╚═╧═╧═══╝} \\
\texttt{```} \\
\texttt{Vous jouez actuellement le tour 7 sur un maximum de 20.} \\
\texttt{Vous avez collectivement accumulé 12 pénalités sur 12.} \\
\texttt{L'autre joueur a envoyé ce message :} \\
\\ 
\\ 
\texttt{"terminé ?"} \\
\\ 
\\ 
\texttt{Quelle est votre prochaine commande ?} \\
            }
        }
    }
    & \\ \\

    \theutterance \stepcounter{utterance}  
    & \multicolumn{4}{p{0.6\linewidth}}{
        \cellcolor[rgb]{0.9,0.9,0.9}{
            \makecell[{{p{\linewidth}}}]{
                \texttt{\tiny{[P1$\rangle$GM]}}
                \texttt{DIRE: terminé !} \\
            }
        }
    }
    & & \\ \\

    \theutterance \stepcounter{utterance}  
    & & & \multicolumn{2}{p{0.3\linewidth}}{
        \cellcolor[rgb]{0.9,0.9,0.9}{
            \makecell[{{p{\linewidth}}}]{
                \texttt{\tiny{[GM$|$GM]}}
                \texttt{DIRE: terminé !} \\
            }
        }
    }
    & & \\ \\

    \theutterance \stepcounter{utterance}  
    & & \multicolumn{4}{p{0.6\linewidth}}{
        \cellcolor[rgb]{0.9,0.9,0.9}{
            \makecell[{{p{\linewidth}}}]{
                \texttt{\tiny{[GM$\rangle$P2]}}
                \texttt{Votre message a été transmis à l'autre joueur.} \\
\\ 
\\ 
\texttt{Vous jouez actuellement le tour 8 sur un maximum de 20.} \\
\texttt{Vous avez collectivement accumulé 12 pénalités sur 12.} \\
\texttt{L'autre joueur a envoyé ce message :} \\
\\ 
\\ 
\texttt{"terminé !"} \\
\\ 
\\ 
\texttt{Quelle est votre prochaine commande ?} \\
            }
        }
    }
    & \\ \\

    \theutterance \stepcounter{utterance}  
    & & & \multicolumn{4}{p{0.6\linewidth}}{
        \cellcolor[rgb]{0.9,0.9,0.9}{
            \makecell[{{p{\linewidth}}}]{
                \texttt{\tiny{[GM$\langle$P2]}}
                \texttt{DIRE: terminé !} \\
            }
        }
    }
     \\ \\

    \theutterance \stepcounter{utterance}  
    & & & \multicolumn{2}{p{0.3\linewidth}}{
        \cellcolor[rgb]{0.9,0.9,0.9}{
            \makecell[{{p{\linewidth}}}]{
                \texttt{\tiny{[GM$|$GM]}}
                \texttt{DIRE: terminé !} \\
            }
        }
    }
    & & \\ \\

    \theutterance \stepcounter{utterance}  
    & & & \multicolumn{2}{p{0.3\linewidth}}{
        \cellcolor[rgb]{0.9,0.9,0.9}{
            \makecell[{{p{\linewidth}}}]{
                \texttt{\tiny{[GM$|$GM]}}
                \texttt{true} \\
            }
        }
    }
    & & \\ \\

    \theutterance \stepcounter{utterance}  
    & & & \multicolumn{2}{p{0.3\linewidth}}{
        \cellcolor[rgb]{0.9,0.9,0.9}{
            \makecell[{{p{\linewidth}}}]{
                \texttt{\tiny{[GM$|$GM]}}
                \texttt{Initial states:} \\
\texttt{```} \\
\texttt{Player 1:} \\
\texttt{ 1234567} \\
\texttt{╔═══╤═╤═╗} \\
\texttt{║C◌◌│◌│◌║ 1} \\
\texttt{║◌I◌│◌├─╢ 2} \\
\texttt{║◌◌◌├─┘◌║ 3} \\
\texttt{║◌◌◌│◌◌◌║ 4} \\
\texttt{╟─┬─┼───╢ 5} \\
\texttt{║◌│◌│◌WT║ 6} \\
\texttt{║◌│◌│◌◌H║ 7} \\
\texttt{╚═╧═╧═══╝} \\
\\ 
\\ 
\texttt{Player 2:} \\
\texttt{ 1234567} \\
\texttt{╔═══╤═╤═╗} \\
\texttt{║◌◌◌│C│◌║ 1} \\
\texttt{║◌W◌│◌├─╢ 2} \\
\texttt{║◌◌◌├─┘◌║ 3} \\
\texttt{║◌◌◌│◌◌H║ 4} \\
\texttt{╟─┬─┼───╢ 5} \\
\texttt{║T│◌│I◌◌║ 6} \\
\texttt{║◌│◌│◌◌◌║ 7} \\
\texttt{╚═╧═╧═══╝} \\
\\ 
\texttt{```} \\
            }
        }
    }
    & & \\ \\

    \theutterance \stepcounter{utterance}  
    & & & \multicolumn{2}{p{0.3\linewidth}}{
        \cellcolor[rgb]{0.9,0.9,0.9}{
            \makecell[{{p{\linewidth}}}]{
                \texttt{\tiny{[GM$|$GM]}}
                \texttt{End states:} \\
\texttt{```} \\
\texttt{Player 1:} \\
\texttt{ 1234567} \\
\texttt{╔═══╤═╤═╗} \\
\texttt{║C◌◌│◌│◌║ 1} \\
\texttt{║◌◌◌│◌├─╢ 2} \\
\texttt{║◌◌T├─┘◌║ 3} \\
\texttt{║◌◌◌│I◌◌║ 4} \\
\texttt{╟─┬─┼───╢ 5} \\
\texttt{║◌│◌│◌W◌║ 6} \\
\texttt{║◌│◌│◌◌H║ 7} \\
\texttt{╚═╧═╧═══╝} \\
\\ 
\\ 
\texttt{Player 2:} \\
\texttt{ 1234567} \\
\texttt{╔═══╤═╤═╗} \\
\texttt{║◌◌◌│C│◌║ 1} \\
\texttt{║◌W◌│◌├─╢ 2} \\
\texttt{║◌◌T├─┘◌║ 3} \\
\texttt{║◌◌◌│I◌◌║ 4} \\
\texttt{╟─┬─┼───╢ 5} \\
\texttt{║◌│◌│H◌◌║ 6} \\
\texttt{║◌│◌│◌◌◌║ 7} \\
\texttt{╚═╧═╧═══╝} \\
\\ 
\texttt{```} \\
            }
        }
    }
    & & \\ \\

    \theutterance \stepcounter{utterance}  
    & & & \multicolumn{2}{p{0.3\linewidth}}{
        \cellcolor[rgb]{0.9,0.9,0.9}{
            \makecell[{{p{\linewidth}}}]{
                \texttt{\tiny{[GM$|$GM]}}
                \texttt{* success: True} \\
\texttt{* lose: False} \\
\texttt{* aborted: False} \\
\texttt{{-}{-}{-}{-}{-}{-}{-}} \\
\texttt{* Shifts: 3.00} \\
\texttt{* Max Shifts: 8.00} \\
\texttt{* Min Shifts: 4.00} \\
\texttt{* End Distance Sum: 11.89} \\
\texttt{* Init Distance Sum: 23.66} \\
\texttt{* Expected Distance Sum: 20.95} \\
\texttt{* Penalties: 12.00} \\
\texttt{* Max Penalties: 12.00} \\
\texttt{* Rounds: 8.00} \\
\texttt{* Max Rounds: 20.00} \\
\texttt{* Object Count: 5.00} \\
            }
        }
    }
    & & \\ \\

\end{supertabular}
}

\end{document}
